\documentclass[12pt]{article}
\usepackage[margin=1in]{geometry}
\usepackage{amsmath,amssymb}
\usepackage{setspace}
\usepackage{hyperref}
\usepackage{natbib}
\usepackage{booktabs}
\usepackage{graphicx}
\usepackage{url}

\onehalfspacing

\title{Jay Gould's Gold Corner of 1869 and the ``Open Market Operation'' Defense:\\
A Sticky-Price Argument Avant la Lettre}

\author{Thomas J.\ Sargent \and Claude}

\date{March 2026 \\ \medskip \small Draft for \emph{Fiscal History} project}

\begin{document}
\maketitle

\begin{abstract}
In September 1869, Jay Gould and James Fisk attempted to corner the gold market 
in New York. When called before the congressional committee chaired by 
James A.\ Garfield, Gould defended his speculative campaign as an ``open 
market operation'' designed to depreciate the greenback, thereby making 
American agricultural exports cheaper in world markets denominated in gold 
and sterling. This paper examines the economic logic of Gould's argument, 
which rested on a sticky-price assumption that domestic greenback prices of 
agricultural goods would not immediately adjust upward in response to a 
rise in the gold premium. We evaluate this defense using the monetary 
economics of the greenback era and connect it to modern theories of 
exchange rate pass-through. We conclude that while the sticky-price 
mechanism Gould invoked has genuine theoretical content, the evidence 
strongly suggests the argument was a post-hoc rationalization crafted 
to secure political cover for a speculative scheme.
\end{abstract}

\tableofcontents

\newpage

%% ============================================================
\section{Introduction}
%% ============================================================

A small date correction is in order at the outset: the events in question 
took place in September 1869, not 1871.\footnote{The Garfield Committee 
report was issued on March 1, 1870, and Henry Adams's article ``The New 
York Gold Conspiracy'' was published in 1870 in the \emph{Westminster 
Review} and subsequently reprinted in \emph{Chapters of Erie} (1871), 
which may be the source of the confusion with the later date.}

Between the spring and fall of 1869, Jay Gould---a director of the Erie 
Railroad---conceived and executed one of the most audacious speculative 
schemes in American financial history: an attempt to corner the market for 
gold on the New York Gold Exchange. The scheme culminated in ``Black 
Friday,'' September 24, 1869, when President Grant ordered Secretary of 
the Treasury George S.\ Boutwell to sell \$4 million in Treasury gold, 
breaking the corner and triggering a panic that sent the gold price 
plummeting from \$160 to \$138 in minutes.

Two important accounts of the episode survive:
\begin{enumerate}
    \item Henry Adams, ``The New York Gold Conspiracy'' (1870), a 
    masterpiece of investigative journalism that appeared in the 
    \emph{Westminster Review} and was reprinted in Adams and Charles 
    Francis Adams, Jr., \emph{Chapters of Erie and Other Essays} (1871).
    \item The \emph{Investigation into the Causes of the Gold Panic: 
    Report of the Majority of the Committee on Banking and Currency} 
    (March 1, 1870), chaired by Representative James A.\ Garfield 
    (41st Congress, 2nd Session, House Report No.\ 31).
\end{enumerate}

In his congressional testimony, Gould offered a remarkable defense of his 
gold purchases. He maintained that his operations were not speculative 
manipulation but rather a form of what we would today call an ``open 
market operation''---an intervention in asset markets designed to raise the 
gold premium (equivalently, to depreciate the greenback against gold), 
which would cheapen American agricultural exports in world markets and 
thereby stimulate the movement of western crops to eastern ports, 
generating freight revenue for his recently acquired Erie Railroad.

This argument is worth examining carefully because it rests on an economic 
assumption---sticky prices---that played a central role in monetary 
debates of the nineteenth century and remains at the heart of modern 
macroeconomics.

%% ============================================================
\section{Institutional Background: The Greenback Era}
%% ============================================================

\subsection{The Dual Monetary Standard}

During and after the Civil War, the United States operated under a 
\emph{dual monetary standard}. The Legal Tender Acts of 1862 and 1863 
authorized the issuance of paper ``greenbacks'' (United States Notes) 
that were legal tender for most domestic transactions but were \emph{not} 
redeemable in gold. Gold continued to circulate and was required for 
payment of customs duties and interest on the federal debt.

The result was a two-price system. Goods could be priced in greenbacks or 
in gold, and the \emph{gold premium}---the greenback price of a gold 
dollar---fluctuated daily on the New York Gold Exchange, established in 
1862. If the gold premium stood at \$135, it meant that \$135 in 
greenbacks were needed to buy \$100 in gold coin. Equivalently, the 
greenback was worth roughly $100/135 \approx 74$ cents in gold.

\subsection{International Trade and the Gold Premium}

International trade was conducted in gold or sterling. American exporters 
who sold wheat in Liverpool received payment in sterling, which could be 
converted into gold at the fixed mint parity. An American farmer, 
however, incurred costs---labor, land, seed, transportation---denominated 
in greenbacks. The profitability of exporting thus depended critically on 
the ratio of the gold (world) price of the commodity to its greenback 
(domestic) cost.

This created a direct link between the gold premium and the incentive to 
export. \emph{If} domestic greenback prices of agricultural goods were 
\emph{slow to adjust} to changes in the gold premium, then a rise in the 
premium (a depreciation of the greenback) would temporarily make American 
exports more profitable. This is the kernel of Gould's argument.

\subsection{Treasury Gold Sales and the Sinking Fund}

Secretary Boutwell pursued an active policy of selling gold from the 
Treasury's reserve of approximately \$100 million in gold bars, using the 
proceeds to retire wartime bonds. These regular gold sales---roughly 
\$1--2 million per week---acted to keep the gold premium low and the 
greenback relatively strong. By September 1, 1869, Boutwell had reduced 
the national debt by \$50 million through these operations.

This Treasury policy was critical to Gould's scheme. If Gould could 
persuade the Grant administration to \emph{suspend} Treasury gold sales, 
the gold premium would naturally rise as greenback liquidity expanded 
relative to the gold supply. Gould used his argument about agricultural 
exports to lobby for exactly this suspension.

%% ============================================================
\section{The Actors}
%% ============================================================

\subsection{Jay Gould}

Jay Gould (1836--1892) had acquired effective control of the Erie Railroad 
by 1868, in partnership with James Fisk and with the aid of the notorious 
``Erie War'' against Cornelius Vanderbilt. The Erie ran from New York City 
to the Great Lakes and was one of the principal trunk lines for carrying 
western agricultural produce to eastern seaports. Gould thus had a 
plausible \emph{commercial} interest in the volume of agricultural 
exports, quite apart from his speculative interest in the gold price. This 
dual interest is what made his defense superficially credible.

\subsection{James Fisk, Jr.}

James Fisk (1835--1872), known as ``Jubilee Jim,'' was Gould's partner in 
the Erie Railroad and joined the gold conspiracy in force during the 
summer of 1869. Fisk was a flamboyant showman whose temperament 
complemented Gould's secretive calculation. Fisk ultimately led the 
public buying campaign on Black Friday while Gould quietly sold.

\subsection{Abel Corbin}

Abel Corbin was a financier of ``shady'' reputation who had married 
Virginia (Jennie) Grant, the younger sister of President Ulysses S.\ 
Grant. Corbin served as the intermediary through whom Gould gained social 
access to the President and attempted to influence Treasury gold policy. 
Gould and Fisk used Corbin's dinners and social occasions to present to 
Grant their argument that high gold prices would help western farmers.

\subsection{President Grant and Secretary Boutwell}

President Grant was initially sympathetic to the argument that Treasury 
gold sales were hurting farmers. In a letter to Boutwell, Grant expressed 
concern that driving the price of gold down would harm agricultural 
interests---a notion, the Garfield Committee concluded, ``planted by 
Gould and Fisk.'' Secretary Boutwell, however, was skeptical. He ``saw 
little merit in either Grant's or Gould's arguments, feeling that the 
government had no place in manipulating the market regardless of who 
benefited.'' Nevertheless, in early September, Boutwell halted the 
non-routine Treasury gold sales at Grant's urging, giving Gould what he 
believed was a ``green light'' to escalate his gold purchases.

Grant reversed course decisively when he learned, around September 
22--23, that Corbin had financial interests tied to the gold speculation. 
On September 24 he ordered Boutwell to sell \$4 million in gold, breaking 
the corner.

%% ============================================================
\section{Gould's Argument: A Reconstruction}
%% ============================================================

Let us reconstruct Gould's economic argument in slightly more formal 
terms than he deployed before the Garfield Committee.

\subsection{The Price Mechanism}

Let $p_g$ denote the gold (world) price of wheat per bushel, expressed in 
gold dollars. In a world of integrated commodity markets (Liverpool--New 
York arbitrage), $p_g$ was approximately given by the Liverpool 
price.\footnote{Transport costs aside, which were themselves paid in 
greenbacks on the domestic leg and in sterling on the ocean leg.}

Let $p_{\$}$ denote the greenback price of wheat in domestic markets 
(e.g., Chicago, New York).

Let $e$ denote the gold premium, i.e., the greenback price of one gold 
dollar. In September 1869, $e$ fluctuated between roughly 130 and 162.

\textbf{Purchasing Power Parity (PPP)} would imply:
\begin{equation}\label{eq:ppp}
    p_{\$} = e \cdot p_g
\end{equation}
That is, the greenback price of wheat should equal the gold price times 
the gold premium. If PPP held continuously, changes in the gold premium 
would be fully and immediately reflected in domestic greenback prices, and 
a depreciation of the greenback would confer \emph{no} competitive 
advantage to exporters.

\subsection{The Sticky-Price Assumption}

Gould's argument implicitly assumed that equation (\ref{eq:ppp}) did 
\emph{not} hold in the short run---that $p_{\$}$ was ``sticky'' and 
adjusted only slowly to changes in $e$. Under sticky domestic prices, an 
increase in $e$ (a depreciation of the greenback against gold) would 
create a wedge:
\begin{equation}\label{eq:wedge}
    \frac{p_g}{p_{\$}/e} = \frac{e \cdot p_g}{p_{\$}} > 1
\end{equation}
whenever $p_{\$}$ failed to keep pace with the rise in $e \cdot p_g$. 
This wedge would make it profitable for merchants to buy American wheat 
at the (unchanged) greenback price, convert greenbacks to gold, and sell 
abroad at the (unchanged) world gold price---earning a margin from the 
greenback's depreciation.

\subsection{The Erie Railroad Connection}

The key link to the Erie Railroad was this: if a rise in the gold premium 
stimulated agricultural exports, the additional grain would flow from the 
western states through trunk-line railroads---including the Erie---to the 
eastern seaboard for shipment to Europe. More freight meant more revenue 
for the Erie.

Thus, Gould argued, his gold purchases served a dual public purpose: 
(a)~they helped western farmers by raising the gold premium and making 
their crops more competitive abroad, and (b)~they stimulated economic 
activity in transportation and commerce.

%% ============================================================
\section{Henry Adams's Account}
%% ============================================================

Henry Adams's ``The New York Gold Conspiracy'' remains the classic 
narrative of the episode. Writing with the controlled fury of a patrician 
moralist, Adams portrayed Gould as a ``spider'' at the center of a web 
of corruption, and Fisk as a swaggering, amoral accomplice.

Adams was particularly scornful of Gould's agricultural-export defense. 
He noted that Gould invoked the argument selectively---deploying it when 
courting political support from Grant and Corbin, while simultaneously 
pursuing aggressive speculative positions that had nothing to do with 
facilitating trade. Adams observed that Gould's gold purchases vastly 
exceeded anything that could be justified by a desire to move the gold 
premium by a few points; by mid-September 1869, Gould and Fisk held 
between \$50 and \$60 million in gold contracts, roughly three times the 
publicly available supply in New York.

Adams also noted the \emph{timing}: Gould began his heaviest purchases 
precisely when he believed Treasury gold sales had been suspended---i.e., 
when his lobbying had created an artificial scarcity. This was 
manipulation of the supply side, not an ``open market operation'' 
designed to achieve a smooth depreciation for the benefit of farmers.

Perhaps most damning was Gould's behavior on Black Friday itself. While 
Fisk was publicly buying gold and driving the price toward \$160, Gould 
was \emph{secretly selling}---liquidating his long positions before the 
inevitable crash. An operator motivated by the public interest of 
promoting agricultural exports would have had no reason for such 
deception.

%% ============================================================
\section{The Garfield Committee Investigation}
%% ============================================================

The Committee on Banking and Currency, chaired by James A.\ Garfield 
(future President), conducted extensive hearings in early 1870. The 
committee's report, issued March 1, 1870, constitutes the most detailed 
documentary record of the episode.

Gould's testimony before the committee was characteristically slippery. 
He maintained that his gold operations were motivated by a desire to help 
agricultural exports and that he had communicated this commercial 
rationale to President Grant. He praised Grant as ``a very pure, 
high-minded man.''

The committee was not persuaded. The report excoriated Gould for his 
manipulation of the gold market, implicated Corbin for exploiting his 
personal connection to Grant, and concluded that Butterfield (the 
assistant treasurer in New York) had served as a ``double agent'' 
providing inside information to the conspirators. The committee cleared 
Grant of personal wrongdoing, though questions lingered about whether 
Julia Grant had received \$25,000 from the speculation.

Garfield and his colleagues understood the economic argument about 
exports and the gold premium. The committee's skepticism was directed not 
at the theoretical mechanism but at \emph{Gould's sincerity in invoking 
it}. The scale, timing, and manner of the gold purchases were 
inconsistent with a benign effort to promote trade; they were consistent 
with a speculative corner designed to extract profits from short-sellers 
and gold borrowers.

%% ============================================================
\section{Evaluating the Sticky-Price Argument}
%% ============================================================

\subsection{Was the Sticky-Price Assumption Reasonable?}

For Gould's argument to have economic merit, the greenback prices of 
agricultural goods would need to have been sticky in the short run---slow 
to adjust to fluctuations in the gold premium.

There is some evidence for this, at least over very short horizons:
\begin{itemize}
    \item \textbf{Contractual rigidities.} Farm wages, land rents, and 
    many input costs in the late 1860s were set in greenbacks by 
    contract and adjusted only periodically.
    \item \textbf{Information lags.} In 1869, information about gold 
    premium fluctuations traveled by telegraph to the major cities but 
    reached rural markets only with delay. The gold premium could move 
    sharply in a single day while local grain prices were set at weekly 
    markets.
    \item \textbf{Transport costs as a buffer.} The cost of 
    transporting grain to the seaboard was denominated in greenbacks 
    and constituted a large fraction of the final price. This 
    greenback-denominated component would not respond immediately to 
    the gold premium.
\end{itemize}

However, the degree and duration of stickiness are crucial. 
Wesley Clair Mitchell's classic study \emph{A History of the Greenbacks} 
(1903) documented that gold-premium fluctuations were transmitted to 
commodity prices with lags of days to weeks, not months. For the 
sustained export stimulus that Gould claimed to be engineering, one 
would need stickiness lasting at least several weeks---a plausible but 
not overwhelming claim.

\subsection{What Does Modern Exchange Rate Economics Say?}

Gould's argument anticipates several ideas in modern international 
macroeconomics:

\begin{enumerate}
    \item \textbf{Incomplete exchange rate pass-through.} A large 
    empirical literature (surveyed by Goldberg and Knetter, 1997, and 
    Burstein and Gopinath, 2014) documents that exchange rate changes 
    are not fully or immediately passed through to domestic prices. The 
    degree of pass-through varies by sector, country, and time horizon, 
    but it is typically incomplete in the short run. This is consistent 
    with the assumption underlying Gould's argument.
    
    \item \textbf{The J-curve.} Following a currency depreciation, the 
    trade balance may initially worsen before improving, because 
    quantities respond more slowly than prices. In Gould's case, even 
    if the gold premium rose, it would take time for additional grain 
    shipments to be organized and transported to port.
    
    \item \textbf{Dornbusch (1976) overshooting.} In a world of sticky 
    goods prices and flexible asset prices, the nominal exchange rate 
    overshoots its long-run equilibrium in response to monetary 
    disturbances. Gould was, in effect, engineering an ``overshoot'' 
    of the gold premium.
    
    \item \textbf{Expenditure-switching effects.} A real depreciation 
    (nominal depreciation with sticky prices) switches expenditure from 
    foreign to domestic goods, stimulating exports and import-competing 
    industries. This is standard textbook analysis (e.g., 
    Obstfeld and Rogoff, 1996).
\end{enumerate}

So the \emph{mechanism} Gould described---depreciation promoting exports 
via sticky domestic prices---is well-grounded in economic theory. The 
question is whether Gould was actually trying to engineer this outcome 
or merely invoking it as cover.

\subsection{The Scale Problem}

One devastating objection to Gould's defense is the \emph{scale} of 
his operations relative to what the stated objective would require.

If the goal were merely to raise the gold premium by a few percentage 
points to promote agricultural exports, one would not need to accumulate 
gold positions of \$50--60 million---three times the publicly available 
supply. A modest position, combined with public advocacy for a change in 
Treasury gold sales policy, would have sufficed.

The enormous scale of Gould's accumulation was \emph{diagnostic of a 
corner}: the goal was not to achieve a new, slightly higher equilibrium 
gold premium but to \emph{squeeze short-sellers}---market participants 
who had borrowed gold and would be forced to buy it back at ever-higher 
prices from the very man who had cornered the supply. The profits from 
a corner come from the squeeze, not from the level of the price.

\subsection{The Behavioral Evidence}

Several features of Gould's behavior are flatly inconsistent with the 
agricultural-export defense:

\begin{enumerate}
    \item \textbf{Secret selling on Black Friday.} While Fisk was 
    publicly bidding gold up to \$160, Gould was secretly selling. An 
    operator genuinely motivated by promoting exports would have no 
    reason to sell into the very price rise he claimed to be engineering 
    for the public good.
    
    \item \textbf{Bribery of officials.} Gould bribed Daniel 
    Butterfield, the assistant treasurer, with \$10,000 to provide 
    advance notice of Treasury gold sales. He attempted to bribe 
    Horace Porter, Grant's personal secretary, with a \$500,000 gold 
    account. He set up accounts for Corbin and (allegedly) for Julia 
    Grant. These are the actions of a speculator seeking inside 
    information, not a public-spirited businessman advocating for 
    sound exchange rate policy.
    
    \item \textbf{Timing relative to the crop season.} If Gould's 
    concern were truly the autumn crop movement, he would have wanted a 
    \emph{sustained} depreciation lasting through the harvest and 
    shipping season (roughly September through December). Instead, the 
    gold price spiked from \$132 to \$160 in three weeks and then 
    crashed back to \$138 in minutes. Such violent instability would 
    have been deeply harmful to the export trade, as merchants require 
    \emph{predictable} exchange rates to enter into forward contracts 
    and make shipping arrangements. Indeed, the aftermath of Black 
    Friday saw a sharp contraction in trade: the \emph{New York 
    Tribune} reported that ``goods ready for export could not be 
    shipped.''
    
    \item \textbf{Actual commodity price effects.} Wheat prices on the 
    Chicago Board of Trade fell from \$1.40 to \$0.77 a bushel in the 
    months following Black Friday. Corn dropped from \$0.95 to \$0.68. 
    These are the opposite of what Gould's theory predicted: if the 
    gold premium were boosting export demand for American grain, 
    greenback prices should have \emph{risen}, not fallen. The actual 
    commodity price collapse reflected the financial disruption and 
    credit contraction that the gold panic caused---a classic case of 
    collateral damage from speculative excess.
\end{enumerate}

%% ============================================================
\section{The Argument as Political Economy}
%% ============================================================

If Gould's agricultural-export argument was not a sincere description of 
his motives, what function did it serve? The answer lies in \emph{political 
economy}: the argument was a device for securing political protection 
for the gold corner.

\subsection{Lobbying the President}

Gould's strategy required the Treasury to \emph{stop selling gold}. As 
long as Boutwell was releasing \$1--2 million per week from the 
Treasury's gold reserve, no private speculator could drive the gold 
premium significantly above its recent range. The Treasury was simply 
too large a player.

To persuade the administration to halt gold sales, Gould needed an 
argument that would resonate with President Grant---a man of limited 
financial sophistication but genuine concern for the welfare of 
farmers and the country's economic recovery. The 
agricultural-export argument was perfectly designed for this audience. 
It framed the suspension of gold sales not as a favor to speculators 
but as a policy to help western farmers---a constituency that Grant, 
a former farmer himself, cared about.

The argument worked, at least temporarily. Grant wrote to Boutwell 
expressing concern that Treasury gold sales were hurting farmers, and 
Boutwell reluctantly suspended non-routine sales in early September 
1869. This suspension was the essential precondition for the escalation 
of the gold corner.

\subsection{Post-Hoc Rationalization}

The agricultural-export argument also served as Gould's legal defense. 
When called before the Garfield Committee, Gould could not plausibly 
deny his massive gold purchases---the records were too clear. His 
strategy was instead to \emph{recharacterize} those purchases as 
motivated by public-spirited commercial considerations rather than 
speculative greed.

This is a familiar pattern in financial history: speculative operators 
retrospectively construct narratives that make their actions appear 
benign or even beneficial. The sophistication of Gould's economic 
reasoning---and the fact that it contains a genuine kernel of 
truth---makes it a particularly elegant specimen of the genre.

%% ============================================================
\section{Parallels and Ironies}
%% ============================================================

\subsection{Gould as an Unofficial Central Banker}

There is a deep irony in Gould's characterization of his gold purchases 
as an ``open market operation.'' In 1869, the United States had no 
central bank. The Second Bank of the United States had been dissolved 
in 1836, and the Federal Reserve would not be established until 1913. 
The Treasury, under Boutwell, was performing a quasi-central-banking 
function through its gold sales and bond purchases.

Gould was, in effect, arguing that the Treasury's ``monetary policy'' 
(selling gold to keep the greenback strong) was misguided, and that his 
\emph{private} open market operations were a corrective. This is an 
extraordinary claim: that a private speculator, acting in his own 
financial interest, could better calibrate the exchange rate than the 
government. The claim is not obviously wrong in principle---governments 
do sometimes pursue misguided exchange rate policies---but it is richly 
ironic coming from a man whose ``calibration'' produced Black Friday.

\subsection{The Real Bills Doctrine and the Gold Premium}

Gould's argument has an interesting connection to the \emph{real bills 
doctrine} that shaped nineteenth-century monetary thinking. Advocates 
of the real bills doctrine held that the quantity of money should 
accommodate the ``needs of trade.'' Applied to the gold market, this 
doctrine would suggest that the gold premium should be allowed to find 
the level that facilitated the smooth flow of agricultural exports---a 
position that, at least superficially, aligns with Gould's advocacy.

The danger, of course, is the one that critics of the real bills 
doctrine have always identified: if the money supply (or in this case, 
the gold premium) is allowed to adjust to accommodate any demand 
presented as a ``need of trade,'' there is no anchor for the price 
level (or the exchange rate), and the result is instability. 
Black Friday was a vivid illustration.

\subsection{Sticky Prices Then and Now}

The sticky-price assumption that Gould deployed informally in 1869 would 
later become one of the central pillars of Keynesian and New Keynesian 
macroeconomics. The idea that nominal rigidities create scope for 
monetary policy to affect real outcomes---output, employment, and the 
trade balance---is the foundation of modern stabilization policy.

Gould's innovation, such as it was, lay in applying this idea to the 
greenback--gold exchange rate in an era when most orthodox economists 
assumed something closer to continuous purchasing power parity. He was 
wrong about his motives but not entirely wrong about the mechanism.

%% ============================================================
\section{Conclusion}
%% ============================================================

Jay Gould's agricultural-export defense of his 1869 gold corner is a 
fascinating episode in the history of economic ideas. The defense rested 
on a sticky-price assumption: that a rise in the gold premium 
(depreciation of the greenback) would make American agricultural 
exports cheaper in world markets because domestic greenback prices 
would be slow to adjust. This mechanism is theoretically sound and 
anticipates modern theories of incomplete exchange rate pass-through.

However, the evidence overwhelmingly indicates that the argument was a 
political instrument, not a sincere description of Gould's commercial 
motives:

\begin{itemize}
    \item The \emph{scale} of Gould's gold accumulation---\$50--60 
    million, three times the public supply---was wildly disproportionate 
    to any plausible export-promotion objective and diagnostic of a 
    speculative corner.
    
    \item Gould's \emph{secret selling} on Black Friday while his 
    partner Fisk publicly bought is inconsistent with any 
    export-promotion rationale.
    
    \item The \emph{bribery} of Treasury officials for inside 
    information is the signature of speculation, not commerce.
    
    \item The \emph{actual outcome}---a violent price spike, a crash, 
    financial disruption, and falling commodity prices---was the 
    opposite of what a successful export-promotion policy would 
    produce.
    
    \item The argument's primary \emph{function} was to lobby President 
    Grant to suspend Treasury gold sales, the essential precondition 
    for the corner, and to provide a post-hoc defense before the 
    Garfield Committee.
\end{itemize}

Gould's defense is best understood as an early and sophisticated example 
of a recurring pattern in financial history: speculative operators 
wrapping private interests in the language of public policy. The 
sophistication of the economic reasoning makes the deception more, not 
less, noteworthy. One suspects that Gould, a man of undeniable 
intelligence, may actually have understood the sticky-price mechanism he 
invoked. But understanding a mechanism and acting on it in good faith 
are very different things.

As Henry Adams mordantly observed, the episode revealed something 
important about the political economy of the Gilded Age: the capacity 
of private financial power to penetrate and manipulate the institutions 
of democratic government. The sticky-price argument was the intellectual 
lubricant for that penetration.

%% ============================================================
\section*{Bibliographic Note}
%% ============================================================

\noindent The principal primary sources are:

\begin{itemize}
    \item Adams, Henry. ``The New York Gold Conspiracy.'' 
    \emph{Westminster Review}, 1870. Reprinted in Adams, Henry, and 
    Charles Francis Adams, Jr. \emph{Chapters of Erie and Other 
    Essays}. Boston: James R.\ Osgood and Company, 1871.
    
    \item Garfield, James A. \emph{Investigation into the Causes of the 
    Gold Panic: Report of the Majority of the Committee on Banking and 
    Currency}. 41st Congress, 2nd Session, House Report No.\ 31. 
    Washington: Government Printing Office, March 1, 1870.
\end{itemize}

\noindent Important secondary sources include:

\begin{itemize}
    \item Ackerman, Kenneth D. \emph{The Gold Ring: Jim Fisk, Jay Gould, 
    and Black Friday, 1869}. New York: Dodd, Mead \& Co., 1988; 
    reprinted 2005, 2011.
    
    \item Klein, Maury. \emph{The Life and Legend of Jay Gould}. 
    Baltimore: Johns Hopkins University Press, 1986.
    
    \item Grodinsky, Julius. \emph{Jay Gould: His Business Career, 
    1867--1892}. Philadelphia: University of Pennsylvania Press, 1957.
    
    \item Mitchell, Wesley Clair. \emph{A History of the Greenbacks}. 
    Chicago: University of Chicago Press, 1903.
    
    \item Mitchell, Wesley Clair. \emph{Gold, Prices, and Wages under 
    the Greenback Standard}. Berkeley: University of California Press, 
    1908.
    
    \item Unger, Irwin. \emph{The Greenback Era: A Social and Political 
    History of American Finance, 1865--1879}. Princeton: Princeton 
    University Press, 1964.
    
    \item White, Ronald C. \emph{American Ulysses: A Life of Ulysses 
    S.\ Grant}. New York: Random House, 2016.
    
    \item Chernow, Ron. \emph{Grant}. New York: Penguin Press, 2017.
\end{itemize}

\noindent On exchange rate pass-through and sticky prices:

\begin{itemize}
    \item Dornbusch, Rudiger. ``Expectations and Exchange Rate 
    Dynamics.'' \emph{Journal of Political Economy} 84, no.\ 6 (1976): 
    1161--1176.
    
    \item Goldberg, Pinelopi K., and Michael M.\ Knetter. ``Goods 
    Prices and Exchange Rates: What Have We Learned?'' \emph{Journal of 
    Economic Literature} 35, no.\ 3 (1997): 1243--1272.
    
    \item Burstein, Ariel, and Gita Gopinath. ``International Prices 
    and Exchange Rates.'' In \emph{Handbook of International Economics}, 
    Vol.\ 4, edited by Gita Gopinath, Elbanan Helpman, and Kenneth 
    Rogoff, 391--451. Amsterdam: Elsevier, 2014.
    
    \item Obstfeld, Maurice, and Kenneth Rogoff. \emph{Foundations of 
    International Macroeconomics}. Cambridge, MA: MIT Press, 1996.
\end{itemize}

\end{document}
